% Source: Weinberg GR, physical PDF pp. 442 and 444 (printed p. 422,
% repeated in the local scan).
% Coverage: Figure 14.2, angular diameter and proper-motion geometry.
% Status: vector reconstruction; source-reviewed and compile-clean.
% Geometry: the two edge rays reach the origin on a cone of half-angle
%           delta/2, as the accompanying text states, so the marked angle
%           delta spans both rays.  Here delta/2 = 8 degrees and the source
%           lies at radius 6.4, putting its edges at (-+0.8907, 6.3378).  The
%           proper diameter D is an arc length across the source, so it is
%           drawn as a circular arc of that radius about the origin, with the
%           straight measuring arrow set above it.
% Uncertainties: source caption states that the angle is deliberately
%                exaggerated.

\begingroup
\renewcommand{\thefigure}{14.2}
\begin{figure}[htbp]
\centering
\begin{tikzpicture}[x=1cm,y=1cm,>=Latex]
  \coordinate (O) at (0,0);
  \coordinate (L) at (-0.8907,6.3378);
  \coordinate (R) at (0.8907,6.3378);

  % Edge rays, travelling from the source down to the observer.
  \draw[line width=0.7pt,
    postaction={decorate},
    decoration={markings,mark=at position 0.45 with {\arrow{Latex}}}]
    (L) -- (O);
  \draw[line width=0.7pt,
    postaction={decorate},
    decoration={markings,mark=at position 0.45 with {\arrow{Latex}}}]
    (R) -- (O);

  % The source itself: proper diameter D, an arc length at radius r_1.
  \draw[line width=0.7pt] (98:6.4) arc[start angle=98,end angle=82,radius=6.4];
  \draw[line width=0.5pt] (0,6.4) -- (O);

  \draw[|<->|] (-0.8907,6.95) -- node[above,inner sep=2pt] {\(D\)}
               (0.8907,6.95);

  % The angular diameter subtended at the observer.
  \draw[<->] (-0.3232,2.30)
    -- node[above,inner sep=1.5pt,fill=white] {\(\delta\)} (0.3232,2.30);

  \node[right,inner sep=4pt] at (O) {\(r=0\)};
  \node[right,inner sep=4pt] at (R) {\(r=r_1\)};
\end{tikzpicture}
\caption{Quantities used in the calculation of angular diameters and proper
motions. The angle \(\delta\) is greatly exaggerated.}
\label{fig:14.2}
\end{figure}
\endgroup
